% Paper 2, §4 — Optimization
% Anchors: kb/notes/training/optimization.md
%          kb/excerpts/muon-moonlight2025.md
%            (#sec-1-challenges, #sec-2-1-update-rule, #sec-2-1-ns,
%             #sec-2-1-norm, #sec-2-2-wd, #sec-2-2-rms, #sec-2-2-hybrid,
%             #sec-2-3-distributed, #sec-3-2-scaling, #sec-3-3-moonlight)

\section{Optimization}
\label{sec:optimization}

The scaling-law section fixed the budget split $(N, D)$ at fixed $C$
and the hyperparameter-transfer protocol that makes any chosen point
trainable; it left the optimizer that spends the FLOPs as a free
variable. Optimizer choice is one of the largest remaining levers on
training compute. \glsterm{adamw}{AdamW} is the production default at every disclosed
frontier lab, with a particular empirical signature
--- update-RMS in the $0.2$--$0.4$ range per matrix parameter
\citep[\S 2.2]{muon-moonlight2025}\footnote{KB excerpt:
\texttt{kb/excerpts/\glsterm{muon}{muon}-moonlight2025.md\#sec-2-2-rms}.} --- that the
2024--2026 alternatives all measure themselves against. \glsterm{lion}{Lion}'s
sign-momentum and \glsterm{sophia}{Sophia}'s diagonal-Hessian preconditioning each
reported $\sim$1.5--2$\times$ training-loss speedups; both gaps largely
closed on downstream evaluation.\deepencite{lion2023 \S 4: head-to-head
LLM training-loss speedup vs.\ AdamW that did not transfer to
downstream-task generalization.} \deepencite{sophia2023 \S 5:
diagonal-Hessian preconditioning, $\sim 2\times$ AdamW training-loss
speedup, downstream-evaluation gap.} \glsterm{per-shape-rms-scaling}{Muon} \citep{muon-moonlight2025},
which orthogonalizes momentum on matrix \glsterm{parameters}{parameters} via a Newton--Schulz
polynomial iteration, is the only published optimizer with a
$\sim\!2\times$ compute-efficiency lift verified at \glsterm{llm}{LLM} scale: the
Moonlight 3B-active / 16B-MoE total run on $5.7$T tokens
\citep[\S 3.3]{muon-moonlight2025}.\footnote{KB excerpt:
\texttt{kb/excerpts/\glsterm{muon}{muon}-moonlight2025.md\#sec-3-3-moonlight}.} The
schedule axis runs in parallel: cosine decay from GPT-3 / Llama-1/2/3
through to the \glsterm{warmup}{Warmup}--Stable--Decay (WSD) family that originated with
MiniCPM \citep{phi4} and is the schedule used by the Moonlight \glsterm{per-shape-rms-scaling}{Muon}
demonstration. Schedule-Free \citep[ICML
2024]{muon-moonlight2025}\deepencite{schedule-free-defazio-2024 \S 1:
removes LR schedule entirely via online primal averaging; limited
LLM-scale validation.} sits as the open frontier.

The thesis of this section is that for matrix-shaped weights ---
which dominate parameter count in a \glsterm{transformer}{Transformer} --- the geometry of
the update is now a first-class design variable. \glsterm{adamw}{AdamW} updates each
coordinate independently under a per-coordinate variance estimate;
\glsterm{per-shape-rms-scaling}{Muon} updates each matrix as a matrix, projecting the momentum onto
its polar factor so all singular directions of the update have equal
energy. The two families' published scaling laws fit the same
\citet{kaplan2020}-style functional form (\cref{eq:6nd}) with
near-identical exponents but different prefactors:
\begin{align}
L_{\mathrm{Muon}}(C)  &\;=\; 2.506 \cdot C^{-0.052},
   \label{eq:muon-scaling-fit} \\
L_{\mathrm{AdamW}}(C) &\;=\; 2.608 \cdot C^{-0.054},
   \label{eq:adamw-scaling-fit}
\end{align}
verbatim from \citet[\S 3.2, Table 3]{muon-moonlight2025}\footnote{KB
excerpt: \texttt{kb/excerpts/\glsterm{muon}{muon}-moonlight2025.md\#sec-3-2-scaling}.}
along the compute-optimal Chinchilla frontier at the Moonlight sweep's
scale. The numerical content of the $2\times$ claim is exactly that
\cref{eq:muon-scaling-fit} sits below \cref{eq:adamw-scaling-fit} by a
near-constant factor, so matched-loss compute requires
$\approx 0.52\,C_{\mathrm{AdamW}}$ \citep[Figure
1a]{muon-moonlight2025}.

\subsection{The optimization step and its symbols}
\label{sec:opt-setup}

The pre-training objective on a corpus $\mathcal{D}$ of token
sequences $\mathbf{x}$ is the next-token cross-entropy
\begin{equation}
\mathcal{L}(\theta) \;=\; -\,\E_{\mathbf{x} \sim \mathcal{D}}\!\left[
   \sum_{t=1}^{T} \log \pi_\theta(x_t \mid x_{<t}) \right],
\label{eq:opt-pretrain-loss}
\end{equation}
minimized by stochastic first-order updates of the form
\begin{equation}
\theta_{t+1} \;=\; \theta_t \,-\, \eta_t \cdot U_t(g_t,\,\theta_t),
\label{eq:opt-step}
\end{equation}
with $g_t = \nabla_\theta \mathcal{L}(\theta_t)$ on a mini-batch,
$\eta_t$ the schedule-controlled learning rate, and $U_t$ the update
returned by the optimizer (which may carry state).
Table~\ref{tab:opt-shapes} fixes notation for the rest of the section.

\begin{table}[h]
\centering
\small
\begin{tabular}{lll}
\toprule
Symbol & Type & Meaning \\
\midrule
$\theta$ & $\R^P$ & full parameter vector ($P \in [10^9, 10^{12}]$ at frontier) \\
$\mathbf{W}$ & $\R^{A \times B}$ & a matrix-shaped parameter (\glsterm{ffn}{FFN}, attention QKV/O) \\
$\mathbf{M}_t$ & $\R^{A \times B}$ & momentum buffer for matrix $\mathbf{W}$ \\
$\mathbf{O}_t$ & $\R^{A \times B}$ & polar-factor approximation $\mathrm{Newton\text{-}Schulz}(\mathbf{M}_t)$ \\
$g_t$ & $\R^P$ & mini-batch gradient at step $t$ \\
$m_t, v_t$ & $\R^P$ & Adam first / second moments \\
$\eta_t$ & $\R$ & schedule-controlled learning rate at step $t$ \\
$\beta_1, \beta_2$ & $\R$ & Adam moment decays (typical $0.9,\,0.95$) \\
$\mu$ & $\R$ & \glsterm{per-shape-rms-scaling}{Muon} momentum decay (typical $0.95$) \\
$\lambda$ & $\R$ & decoupled weight-decay coefficient (typical $0.1$) \\
$\varepsilon$ & $\R$ & numerical-stability constant (typical $10^{-8}$) \\
$C$ & FLOPs & total training compute, $C \approx 6 N D$ from \cref{eq:6nd} \\
\bottomrule
\end{tabular}
\caption{Symbols used in \cref{sec:optimization}. Matrix parameters
$\mathbf{W} \in \R^{A \times B}$ are the only shapes Muon updates;
embeddings, the LM head, and 1-D parameters (RMSNorm gains, biases)
stay on AdamW \citep[\S 2.2]{muon-moonlight2025}.\protect\footnote{KB
excerpt: \texttt{kb/excerpts/muon-moonlight2025.md\#sec-2-2-hybrid}.}}
\label{tab:opt-shapes}
\end{table}

\subsection{AdamW, the production default}
\label{sec:opt-adamw}

Adam \citep[Kingma \&
Ba 2015]{muon-moonlight2025}\deepencite{kingma2015-adam: Adam
optimizer original, $(\beta_1, \beta_2)$ moment definitions, the
$\hat{m}_t / (\sqrt{\hat{v}_t} + \varepsilon)$ update rule used
verbatim below.} keeps two exponential moving averages of the
gradient and its square,
\begin{equation}
m_t \;=\; \beta_1 m_{t-1} + (1 - \beta_1) g_t, \qquad
v_t \;=\; \beta_2 v_{t-1} + (1 - \beta_2) g_t^2,
\label{eq:adam-moments}
\end{equation}
applies bias correction, and uses the ratio
$\hat{m}_t / (\sqrt{\hat{v}_t} + \varepsilon)$ as the per-coordinate
update direction, with $\hat{m}_t = m_t / (1 - \beta_1^t)$ and
$\hat{v}_t = v_t / (1 - \beta_2^t)$. The ``W'' in \glsterm{adamw}{AdamW}
\citep[Loshchilov \&
Hutter 2019]{muon-moonlight2025}\deepencite{loshchilov2019-adamw:
decoupled weight-decay derivation; required citation for the
$- \eta_t \lambda \theta_t$ term being applied \emph{after} Adam's
preconditioning rather than folded into $g_t$.} is decoupled weight
decay: rather than fold $\lambda \theta$ into $g_t$ (which would put
weight decay inside Adam's per-coordinate rescaling), \glsterm{adamw}{AdamW} applies
the decay to $\theta$ directly:
\begin{equation}
\theta_{t+1} \;=\; \theta_t \,-\, \eta_t \!\left[
   \frac{\hat{m}_t}{\sqrt{\hat{v}_t} + \varepsilon} \;+\; \lambda\,\theta_t
\right].
\label{eq:adamw}
\tag{2}
\end{equation}
\Cref{eq:adamw} is the universal \glsterm{adamw}{AdamW} step at \glsterm{llm}{LLM} scale --- LLaMA
\citep{meta-llama3}, \glsterm{precision-pinning}{DeepSeek-V3} \citep[\S 4]{deepseek-v3}, and Phi-4
\citep{phi4} all use it, with typical
$(\beta_1, \beta_2) = (0.9, 0.95)$, $\lambda = 0.1$, and $\varepsilon
= 10^{-8}$.

\paragraph{The empirical update-RMS signature.}
For each matrix parameter $\mathbf{W} \in \R^{A \times B}$, the \glsterm{adamw}{AdamW}
update has root-mean-square magnitude
$\| \eta_t \, \hat{m}_t / (\sqrt{\hat{v}_t} + \varepsilon) \|_2 /
\sqrt{AB}$ that is ``usually around $0.2$ to $0.4$'' \citep[\S
2.2]{muon-moonlight2025}.\footnote{KB excerpt:
\texttt{kb/excerpts/\glsterm{muon}{muon}-moonlight2025.md\#sec-2-2-rms}.} This number
is the calibration target every alternative measures itself against:
matching it lets one reuse \glsterm{adamw}{AdamW}'s tuned $(\eta, \lambda)$ across
optimizers, sparing a frontier-scale sweep. \Cref{sec:opt-muon}
reports how the \glsterm{per-shape-rms-scaling}{Muon} family lands exactly on this number by design.

\paragraph{Optimizer-state cost.}
\glsterm{adamw}{AdamW} carries $2P$ words of moment state $(m, v)$ on top of the
\glsterm{parameters}{parameters} and gradients --- $4P$ words total in mixed-precision
training when an FP32 master copy is kept (see
\cref{sec:mixed-precision}). For a $70$B-parameter model in FP32 the
optimizer-state footprint alone is $\sim 1.1$\,TB
\citep[\S 2]{muon-moonlight2025}.\deepencite{adamw-state-cost: standard
back-of-envelope --- $4P$ words for params + grads + Adam $(m, v)$;
referenced rather than canonically cited because it is folklore in
the distributed-training literature.}

\subsection{Lion and Sophia --- the training-loss-conditional gap}
\label{sec:opt-lion-sophia}

Two \glsterm{adamw}{AdamW} alternatives published in 2023 reported
$\sim$1.5--2$\times$ speedups on training loss and have been
extensively re-evaluated since.

\paragraph{Lion --- sign-momentum.}
\glsterm{lion}{Lion} \citep[Chen et
al.\ 2023]{muon-moonlight2025}\deepencite{lion2023 \S 2: sign-momentum
update rule with single moment buffer; arXiv 2302.06675.} replaces
the Adam ratio with the sign of an exponentially smoothed gradient,
giving
\begin{equation}
c_t \;=\; \beta_1 m_{t-1} + (1 - \beta_1) g_t, \qquad
\theta_{t+1} \;=\; \theta_t \,-\, \eta_t \!\left(
   \mathrm{sign}(c_t) + \lambda \theta_t
\right),
\label{eq:lion-update}
\end{equation}
with the moment buffer updated as
$m_t = \beta_2 m_{t-1} + (1-\beta_2) g_t$. A single moment buffer
halves the optimizer-state memory relative to Adam and the per-step
cost is one elementwise sign rather than a square root. \glsterm{lion}{Lion}'s
headline claim is wall-clock speedup at matched final loss on vision
transformers; for LLMs the gap closed in head-to-head
comparisons.\deepencite{three-optimizer-2025 (arXiv:2507.08472):
controlled LLM head-to-head reporting AdamW catching up or winning
on downstream evaluation; cited as Phase-1-tier signal in the KB
note.}

\paragraph{Sophia --- diagonal-Hessian preconditioning.}
\glsterm{sophia}{Sophia} \citep[Liu, Li et al.\ 2023, ICLR
2024]{muon-moonlight2025}\deepencite{sophia2023 \S 3 (arXiv
2305.14342): stochastic diagonal-Hessian estimator $h_t$, the
clipped second-order update
$\theta_{t+1} = \theta_t - \eta_t \mathrm{clip}(m_t / \max(h_t,
\varepsilon), \rho)$, and the headline ``$2\times$ AdamW training-loss
speedup'' figure.} replaces Adam's gradient-variance preconditioner
$\sqrt{\hat{v}_t}$ with a stochastic diagonal-Hessian estimate
$h_t$, computed via a Hutchinson-style \glsterm{probe-2}{probe} on the network's curvature
operator, and clips the resulting step. The reported $\sim\!2\times$
speedup over \glsterm{adamw}{AdamW} on training loss held in published replications.
The downstream gap is the load-bearing complication: \glsterm{adamw}{AdamW} catches up
or surpasses \glsterm{sophia}{Sophia} on standard evaluation suites at matched
\emph{wall-clock} time, despite \glsterm{sophia}{Sophia}'s lower training loss.

\begin{contradiction}
\label{contra:opt-loss-vs-downstream}
Both \glsterm{lion}{Lion} and \glsterm{sophia}{Sophia} exhibit the same pattern. Lower training loss at
matched compute, but downstream evaluation in the closing-the-gap or
gap-reversed regime once both optimizers run to a comparable budget
and the downstream metric (\glsterm{mmlu}{MMLU}, GSM8K, HumanEval, BBH-class)
replaces validation NLL. The mechanistic interpretation has not been
isolated from the experimental confound (LR sweep, weight-decay
sweep, schedule shape interaction). Treat reported optimizer
speedups as \emph{training-loss-conditional} until verified on the
downstream evaluation panel; this is the operational lesson of the
2024--2025 \glsterm{llm}{LLM} optimizer literature
\citep[\S 1]{muon-moonlight2025}.\protect\footnote{KB excerpt:
\texttt{kb/excerpts/\glsterm{muon}{muon}-moonlight2025.md\#sec-1-challenges}.}
\end{contradiction}

\subsection{Muon --- orthogonalized momentum on matrix parameters}
\label{sec:opt-muon}

\citet{muon-moonlight2025} scale up the \glsterm{per-shape-rms-scaling}{Muon} optimizer
(originally K.\ Jordan et al.\ 2024)\deepencite{jordan2024-muon:
original ``Muon: an optimizer for hidden layers in neural networks''
write-up; the Moonlight paper is the LLM-scale extension that this
section primarily cites.} to \glsterm{llm}{LLM} training. \glsterm{per-shape-rms-scaling}{Muon}'s central
observation is that a matrix parameter $\mathbf{W} \in \R^{A \times B}$
is not a flat vector. The gradient $\nabla \mathcal{L} \in \R^{A
\times B}$ has singular structure --- some directions in row /
column space accumulate signal faster than others --- and the \glsterm{adamw}{AdamW}
update preserves this anisotropy because per-coordinate rescaling is
basis-dependent. \glsterm{per-shape-rms-scaling}{Muon}'s update orthogonalizes the momentum before
the step, replacing $\mathbf{M}_t$ with its polar factor $\mathbf{O}_t
= \mathbf{U} \mathbf{V}^\top$ in the SVD $\mathbf{M}_t = \mathbf{U}
\boldsymbol{\Sigma} \mathbf{V}^\top$. Verbatim from
\citet[\S 2.1, eq.\ 1]{muon-moonlight2025}\footnote{KB excerpt:
\texttt{kb/excerpts/\glsterm{muon}{muon}-moonlight2025.md\#sec-2-1-update-rule}.}:
\begin{equation}
\begin{aligned}
\mathbf{M}_t &\;=\; \mu \, \mathbf{M}_{t-1} \,+\, \nabla \mathcal{L}_t(\mathbf{W}_{t-1}) \\
\mathbf{O}_t &\;=\; \mathrm{Newton\text{-}Schulz}(\mathbf{M}_t) \\
\mathbf{W}_t &\;=\; \mathbf{W}_{t-1} \,-\, \eta_t \,\mathbf{O}_t.
\end{aligned}
\label{eq:muon-vanilla}
\tag{3}
\end{equation}
$\mathbf{O}_t$ approximates the polar factor of $\mathbf{M}_t$, i.e.\
the matrix whose SVD has the same singular vectors as $\mathbf{M}_t$
but whose singular values are all unity:
$(\mathbf{M}_t \mathbf{M}_t^\top)^{-1/2} \mathbf{M}_t = \mathbf{U}
\mathbf{V}^\top$ \citep[\S 2.1]{muon-moonlight2025}. The
orthogonalization ``can ensure that the update matrices are
isomorphic, preventing the weight from learning along a few dominant
directions'' \citep[\S 2.1]{muon-moonlight2025}.

\paragraph{The Newton--Schulz iteration.}
The polar factor of $\mathbf{M}_t$ is computed by a matrix polynomial
iteration that avoids the SVD entirely. Initialize
$\mathbf{X}_0 = \mathbf{M}_t / \|\mathbf{M}_t\|_F$ (Frobenius
normalization keeps the iteration in the convergence basin); then
\begin{equation}
\mathbf{X}_k \;=\; a\,\mathbf{X}_{k-1}
   \;+\; b\,(\mathbf{X}_{k-1} \mathbf{X}_{k-1}^\top)\,\mathbf{X}_{k-1}
   \;+\; c\,(\mathbf{X}_{k-1} \mathbf{X}_{k-1}^\top)^2\,\mathbf{X}_{k-1},
\label{eq:newton-schulz}
\tag{4}
\end{equation}
verbatim from \citet[\S 2.1, eq.\ 2]{muon-moonlight2025}\footnote{KB
excerpt: \texttt{kb/excerpts/\glsterm{muon}{muon}-moonlight2025.md\#sec-2-1-ns}.}
with coefficients
$(a, b, c) = (3.4445,\, -4.7750,\, 2.0315)$
\citep[K.\ Jordan et al.\ 2024, cited in \S
2.1]{muon-moonlight2025}.\footnote{KB excerpt:
\texttt{kb/excerpts/\glsterm{muon}{muon}-moonlight2025.md\#sec-2-1-ns}.} Five
iterations suffice in production (the polynomial converges to the
polar factor on each side of the singular-\glsterm{value}{value} spectrum at a quintic
rate near unity). The output $\mathbf{X}_N$ is $\mathbf{O}_t$ in
\cref{eq:muon-vanilla}.

\begin{analogy}
\label{analogy:polar}
$\mathbf{O}_t$ is the polar-decomposition projection of
$\mathbf{M}_t$ onto the manifold of orthogonal matrices in the
Frobenius geometry. If $\mathbf{M}_t = \mathbf{U} \boldsymbol{\Sigma}
\mathbf{V}^\top$ is the SVD, then $\mathbf{O}_t = \mathbf{U}
\mathbf{V}^\top$ is the closest orthogonal matrix to $\mathbf{M}_t$
under any unitarily-invariant norm
\citep[Mirsky 1960]{muon-moonlight2025}.\deepencite{mirsky1960-symmetric-gauge:
canonical reference for the polar projection being the
unitarily-invariant-norm-closest orthogonal matrix to a given matrix.}
The Newton--Schulz iteration in \cref{eq:newton-schulz} computes this
projection without ever forming the SVD --- a polynomial in
$\mathbf{X}_{k-1} \mathbf{X}_{k-1}^\top$ approximates the inverse
square root $(\mathbf{X} \mathbf{X}^\top)^{-1/2}$, and the
$(a, b, c)$ coefficients are the optimal cubic-in-singular-values fit
that maps the Frobenius-normalized spectrum of $\mathbf{X}_0$ onto
$\{1\}$ in $N$ iterations. Returning to math: the canonical \glsterm{per-shape-rms-scaling}{Muon}
update in \cref{eq:muon-vanilla} reduces to ``take the polar factor
of momentum, scale by $\eta_t$, subtract,'' and \cref{eq:newton-schulz}
is the SVD-free numerical recipe for the polar factor.
\end{analogy}

\paragraph{Two scale-stable ingredients.}
\citet[\S 2.2]{muon-moonlight2025} report that vanilla
\cref{eq:muon-vanilla} converges faster than \glsterm{adamw}{AdamW} at small scale but
the gain diminishes at \glsterm{llm}{LLM} scale because of two failure modes that
the original formulation does not address. The fixes are
load-bearing.

\textbf{Mandatory weight decay.} Without an explicit decay term,
``both the weight and the layer output's RMS keep growing to a large
scale, exceeding the high-precision range of bf16, which might hurt
the model's performance''
\citep[\S 2.2]{muon-moonlight2025}.\footnote{KB excerpt:
\texttt{kb/excerpts/\glsterm{muon}{muon}-moonlight2025.md\#sec-2-2-wd}.} The fix is
to fold \glsterm{adamw}{AdamW}-style decoupled weight decay into the \glsterm{per-shape-rms-scaling}{Muon} step:
\begin{equation}
\mathbf{W}_t \;=\; \mathbf{W}_{t-1} \,-\, \eta_t \!\left(
   \mathbf{O}_t \,+\, \lambda \mathbf{W}_{t-1}
\right),
\label{eq:muon-wd}
\end{equation}
verbatim from \citet[\S 2.2, eq.\ 3]{muon-moonlight2025}.\footnote{KB
excerpt: \texttt{kb/excerpts/\glsterm{muon}{muon}-moonlight2025.md\#sec-2-2-wd}.}
With \cref{eq:muon-wd}, \glsterm{per-shape-rms-scaling}{Muon} outperforms both vanilla \glsterm{per-shape-rms-scaling}{Muon} and \glsterm{adamw}{AdamW}
in the over-train regime that \cref{tab:frontier-NDC} of
\cref{sec:scaling-laws} characterizes.

\textbf{Per-shape update-RMS adjustment.} Lemma 1 of
\citet[\S 2.2]{muon-moonlight2025}\footnote{KB excerpt:
\texttt{kb/excerpts/\glsterm{muon}{muon}-moonlight2025.md\#sec-2-2-rms}.} states that
for a full-rank matrix parameter of shape $A \times B$, the
theoretical \glsterm{per-shape-rms-scaling}{Muon} update RMS is $\sqrt{1/\max(A, B)}$. Tall-thin
matrices (large $\max(A, B)$, e.g.\ \glsterm{ffn}{FFN} $\mathbf{W}_{\text{up}} \in
\R^{4D \times D}$) get under-updated; small per-head matrices
(treating each KV head in \glsterm{gqa}{GQA} \citep{ainslie2023} or \glsterm{multi-head-latent-attention}{MLA}
\citep{deepseek-v2} as a separate parameter) get over-updated and
destabilize training. The fix is to multiply each matrix's update by
$\sqrt{\max(A, B)}$ and a constant factor $0.2$ that lands in \glsterm{adamw}{AdamW}'s
empirical $0.2$--$0.4$ band. The full step is
\begin{equation}
\mathbf{W}_t \;=\; \mathbf{W}_{t-1} \,-\, \eta_t \!\left(
   0.2 \,\mathbf{O}_t \,\sqrt{\max(A, B)} \;+\; \lambda \mathbf{W}_{t-1}
\right),
\label{eq:muon-final}
\tag{4'}
\end{equation}
verbatim from \citet[\S 2.2, eq.\ 4]{muon-moonlight2025}.\footnote{KB
excerpt: \texttt{kb/excerpts/\glsterm{muon}{muon}-moonlight2025.md\#sec-2-2-rms}.}
The $0.2$ factor is exactly what calibrates the update RMS to \glsterm{adamw}{AdamW}'s
typical $0.2$--$0.4$ range, so \glsterm{adamw}{AdamW}'s tuned $(\eta, \lambda)$ ``can
directly \emph{reuse}'' for \glsterm{per-shape-rms-scaling}{Muon} \citep[\S
2.2]{muon-moonlight2025}.\footnote{KB excerpt:
\texttt{kb/excerpts/\glsterm{muon}{muon}-moonlight2025.md\#sec-2-2-rms}.} This is the
non-trivial production property: the optimizer change does not
require a fresh hyperparameter sweep at frontier scale.

\paragraph{Hybrid Muon + AdamW.}
\glsterm{per-shape-rms-scaling}{Muon} updates only matrix \glsterm{parameters}{parameters} \citep[\S
2.2]{muon-moonlight2025}.\footnote{KB excerpt:
\texttt{kb/excerpts/\glsterm{muon}{muon}-moonlight2025.md\#sec-2-2-hybrid}.}
Embeddings ($\R^{V \times D}$ are 2-D but treated as embeddings, not
matrices, in the \glsterm{per-shape-rms-scaling}{Muon} convention because the row dim $V$ is not a
hidden-dim), the LM head, \glsterm{rmsnorm}{RMSNorm} gains, and any other 1-D parameter
stay on \glsterm{adamw}{AdamW}. A production training run with \glsterm{per-shape-rms-scaling}{Muon} is therefore a
\emph{hybrid optimizer}: \glsterm{per-shape-rms-scaling}{Muon} on the body, \glsterm{adamw}{AdamW} on the periphery,
sharing $(\eta_t, \lambda)$ across both. This is exactly the
hybridization Moonlight \citep[\S 3]{muon-moonlight2025} reports.

\begin{intuition}
\label{intuit:muon-vs-adamw}
Adam's per-coordinate $1/\sqrt{\hat{v}_t}$ implicitly assumes
\glsterm{parameters}{parameters} live in a coordinate system the optimizer can read off ---
the natural basis of $\R^P$. \glsterm{per-shape-rms-scaling}{Muon}'s polar projection treats matrix
\glsterm{parameters}{parameters} as matrices: rotations of the row or column basis leave
$\mathbf{O}_t$ equivariant in a way Adam's update is not. Bernstein
et al.\ 2024 \citep[\S 2.1]{muon-moonlight2025}\protect\footnote{KB
excerpt: \texttt{kb/excerpts/\glsterm{muon}{muon}-moonlight2025.md\#sec-2-1-norm}.}
formalize the difference as a norm constraint on the implicit
steepest-descent step: Adam's norm is the ``Max-of-Max'' (
$\ell_\infty$ over coordinates), which bakes in a coordinate basis;
\glsterm{per-shape-rms-scaling}{Muon}'s is a \glsterm{schatten-p-norm}{Schatten-$p$ norm} for some large $p$, a function of
singular values, which is rotation-invariant. Returning to canonical
form: \cref{eq:newton-schulz} produces $\mathbf{O}_t$ with all
singular values equal to $1$ \citep[\S
2.1]{muon-moonlight2025}, so the matrix update has equal energy in
every direction in the row / column subspaces of $\mathbf{M}_t$ ---
no coordinate is implicitly privileged by gradient magnitude, only
by direction.
\end{intuition}

\paragraph{The 52\%-FLOPs scaling-law claim.}
\citet[\S 3.2]{muon-moonlight2025}\footnote{KB excerpt:
\texttt{kb/excerpts/\glsterm{muon}{muon}-moonlight2025.md\#sec-3-2-scaling}.} run a
multi-scale sweep along the Chinchilla compute-optimal frontier and
fit power laws of the form $L = a C^{-\alpha}$ to validation loss
(see \cref{eq:muon-scaling-fit,eq:adamw-scaling-fit}). Solving
$L_{\mathrm{Muon}}(\eta C^*) = L_{\mathrm{AdamW}}(C^*)$ for the
matched-loss compute ratio $\eta = C_{\mathrm{Muon}} /
C_{\mathrm{AdamW}}$ gives $\eta \approx 0.52$
\citep[Figure 1a]{muon-moonlight2025}: ``\glsterm{per-shape-rms-scaling}{Muon} only requires about
52\% training FLOPs to match the performance of \glsterm{adamw}{AdamW} under
compute optimal setting'' \citep[\S
3.2]{muon-moonlight2025}.\footnote{KB excerpt:
\texttt{kb/excerpts/\glsterm{muon}{muon}-moonlight2025.md\#sec-3-2-scaling}.} The
sweep covers up to the Moonlight production scale; whether the
constant factor in $\eta \approx 0.52$ holds at $70$B+ dense or
$600$B+ MoE is untested.

\paragraph{The Moonlight demonstration.}
The largest published \glsterm{per-shape-rms-scaling}{Muon} run is Moonlight, a $2.24$B-active /
$15.29$B-total parameter Mixture-of-Experts (or $3$B-active / $16$B-
total when including embedding) trained on $5.7$T tokens
\citep[\S 3.3]{muon-moonlight2025}.\footnote{KB excerpt:
\texttt{kb/excerpts/\glsterm{muon}{muon}-moonlight2025.md\#sec-3-3-moonlight}.} The
architecture is the \glsterm{precision-pinning}{DeepSeek-V3}-small variant from \citet{deepseek-v3}
with the auxiliary-loss-free MoE load-balancing scheme; the precision
recipe is BF16 working with FP32 \glsterm{master-weights}{master weights}. The schedule is the
WSD form developed in \cref{sec:opt-schedules}.

\subsection{The schedule axis --- cosine, WSD, Schedule-Free}
\label{sec:opt-schedules}

The schedule $\eta_t$ is the second optimizer-side compute lever.
Two schedules dominate frontier \glsterm{llm}{LLM} training, with a third at the
research frontier.

\paragraph{Cosine decay.}
\citep[Loshchilov \&
Hutter 2017]{muon-moonlight2025}\deepencite{loshchilov2017-cosine:
SGDR / cosine annealing with restarts.}
\begin{equation}
\eta_t \;=\; \eta_{\min} + \tfrac{1}{2}\,(\eta_{\max} - \eta_{\min})\,\Bigl(1 + \cos(\pi\, t / T)\Bigr),
\label{eq:cosine}
\end{equation}
preceded by linear \glsterm{warmup}{warmup} over $\sim$2k steps. Used by GPT-3
\citep{brown2020}, Llama-1/2/3 \citep{meta-llama3}, and most pre-2024
runs. The structural cost of \cref{eq:cosine} is that $T$, the total
training-step count, must be committed at \glsterm{warmup}{warmup} --- continuing
training past $T$ requires a fresh schedule with a discontinuity at
the resumption point. \Cref{sec:scaling-lrhorizon} flagged this as the
operational reason behind the Kaplan-vs-Chinchilla LR-horizon bug:
intermediate-loss readouts at $D' \ll T$ are taken at high-LR
mid-training, not at the LR-decay endpoint.

\paragraph{Warmup--Stable--Decay (WSD).}
WSD --- originated in MiniCPM \citep{phi4}\deepencite{minicpm2024
(arXiv:2404.06395) \S 3: WSD origin paper. The Phi-4 cite stands in
because the bib has Phi-4 and not the MiniCPM paper.} with a
theoretical justification arriving in
2024--2026\deepencite{wsd-river-valley-2024 (arXiv:2410.05192) \S 2:
``river-valley'' loss-landscape view that motivates a flat-then-decay
schedule.} \deepencite{wsd-optimal-2026 (arXiv:2602.06797) \S 4: WSD
derived as the optimal schedule under a functional scaling-laws
analysis.} --- decomposes the training run into three phases:
\begin{equation}
\eta_t \;=\;
\begin{cases}
\eta_{\max}\,\dfrac{t}{T_w}, & 0 \leq t \leq T_w \quad\text{(warmup)} \\[6pt]
\eta_{\max}, & T_w < t \leq T_w + T_s \quad\text{(stable)} \\[6pt]
\mathrm{decay}(t - T_w - T_s,\, \eta_{\max} \to \eta_{\min}), & T_w + T_s < t \leq T \quad\text{(cooldown)}
\end{cases}
\label{eq:wsd}
\end{equation}
where the cooldown is typically linear or cosine over the final
$10$--$20\%$ of training. The decisive practical advantage is that
the stable phase has no fixed endpoint --- training can be continued
indefinitely at $\eta_{\max}$ before committing to a cooldown,
making continuous-training and mid-training mixture changes routine.

The Moonlight \glsterm{per-shape-rms-scaling}{Muon} run uses an explicit WSD with a quirk: the
cooldown is preceded by an LR \emph{re-rise} that loads
high-quality data into a controlled mini-\glsterm{warmup}{warmup}-then-decay phase.
Verbatim from \citet[\S
3.3]{muon-moonlight2025}\footnote{KB excerpt:
\texttt{kb/excerpts/\glsterm{muon}{muon}-moonlight2025.md\#sec-3-3-moonlight}.}:
\begin{equation}
\eta_t^{\mathrm{Moonlight}} \;=\;
\begin{cases}
0 \;\to\; 4.2 \times 10^{-4} & \text{linear over $[0,\,33\text{B}]$ tokens (warmup)} \\[4pt]
4.2 \times 10^{-4} \;\to\; 4.2 \times 10^{-5} & \text{cosine over $[33\text{B},\,5.2\text{T}]$ tokens (stable)} \\[4pt]
4.2 \times 10^{-5} \;\to\; 1 \times 10^{-4} & \text{linear in 100 steps (cooldown rise)} \\[4pt]
1 \times 10^{-4} \;\to\; 0 & \text{linear over $500$B tokens (cooldown decay)}
\end{cases}
\label{eq:moonlight-schedule}
\end{equation}
Phase 3+4 is the cooldown stage in which ``we use the highest
quality data, focusing on math, code, and reasoning'' \citep[\S
3.3]{muon-moonlight2025}. This is the same data-injection pattern
that OLMo-2's Dolmino mix uses (\cref{sec:adaptation-and-merging}):
late-training, small-$\eta$ exposure to curated data has outsized
effect on final loss because the small-$\eta$ regime acts like
SGD-on-a-quadratic near a minimum, where averaging dynamics dominate.

\paragraph{Schedule-Free.}
\citep[Defazio et
al.\ 2024]{muon-moonlight2025}\deepencite{schedule-free-defazio-2024
\S 2--3 (ICML 2024): online primal-averaging update with no LR
schedule; bound matches optimally-scheduled SGD/AdamW under standard
convex assumptions; LLM-scale validation pending.} removes the LR
schedule entirely via online primal averaging:
$\theta_{t+1} = \theta_t - \eta\, U_t$ with $\eta$ \emph{constant},
and a separate evaluation iterate
$\bar{\theta}_t = (1 - c_t)\,\bar{\theta}_{t-1} + c_t\,\theta_t$ for
some averaging weight $c_t$. The bound matches optimally-scheduled
SGD / \glsterm{adamw}{AdamW} under convex assumptions; large-scale \glsterm{llm}{LLM} validation is
limited as of 2026. Schedule-Free's appeal at frontier scale is the
same as WSD's --- no committed total-step count --- but with the
additional removal of the cooldown shape decision; the unresolved
question is whether the cooldown-as-data-mixture-vehicle role of WSD
\eqref{eq:moonlight-schedule} is replicable in a schedule-free
setting.

\subsection{Distributed Muon --- a forward link to \S 5}
\label{sec:opt-distributed-muon}

The optimizer-state partitioning that ZeRO-1 / \glsterm{fsdp}{FSDP} exploit for
\glsterm{adamw}{AdamW} \citep[\S 2.2.4]{torchtitan2024} does not directly apply to
\glsterm{per-shape-rms-scaling}{Muon}: the Newton--Schulz iteration in \cref{eq:newton-schulz}
operates on the full $A \times B$ momentum matrix, not on the
per-\glsterm{data-parallel}{DP}-rank shard. \citet[\S
2.3]{muon-moonlight2025}\footnote{KB excerpt:
\texttt{kb/excerpts/\glsterm{muon}{muon}-moonlight2025.md\#sec-2-3-distributed}.}
introduce Distributed \glsterm{per-shape-rms-scaling}{Muon}, which partitions optimizer state on \glsterm{data-parallel}{DP} as
ZeRO-1 does and adds two operations:
\begin{enumerate}
\item \textbf{\glsterm{data-parallel}{DP} gather.} For each matrix parameter, gather the
\glsterm{data-parallel}{DP}-partitioned gradients into the full $A \times B$ matrix on the
local rank.
\item \textbf{Calculate full update.} Run \cref{eq:newton-schulz} on
the full gradient matrix to produce $\mathbf{O}_t$, then discard the
unused rows and apply the local partition's slice of \cref{eq:muon-final}
to update the local parameter shard.
\end{enumerate}
The latency overhead is reported at $1$\%--$3$\% of forward-backward
time \citep[\S 2.3]{muon-moonlight2025}.\footnote{KB excerpt:
\texttt{kb/excerpts/\glsterm{muon}{muon}-moonlight2025.md\#sec-2-3-distributed}.}
\Cref{sec:distributed-training} expands on how each optimizer
composes with the full $G = d \cdot t \cdot c \cdot p \cdot e$
DeviceMesh \citep[\S 4]{torchtitan2024}: how Distributed \glsterm{per-shape-rms-scaling}{Muon}
interacts with \glsterm{tensor-parallel}{TP} / \glsterm{pipeline-parallel}{PP} / \glsterm{expert-parallel}{EP}, why some axes are
optimizer-state-friendly and others are not, and the open question of
whether the $1$\%--$3$\% overhead scales to $32$k+ GPUs.

\subsection{Open frontiers}
\label{sec:opt-open}

\begin{contradiction}
\label{contra:muon-frontier-scale}
The \glsterm{per-shape-rms-scaling}{Muon}-Moonlight \glsterm{scaling-law}{scaling law} \eqref{eq:muon-scaling-fit} fits a
sweep that tops out at the Moonlight $16$B-MoE / $5.7$T-token point.
The matched-loss compute ratio $\eta \approx 0.52$ is a
\emph{constant} extrapolation: the ratio of two power laws with
nearly identical exponents ($-0.052$ vs.\ $-0.054$) stays
approximately constant in $C$, so the prediction is that \glsterm{per-shape-rms-scaling}{Muon}'s lift
holds at $70$B+ dense and $600$B+ MoE. This is also the open
question \citet{muon-moonlight2025} raise themselves \citep[\S
1]{muon-moonlight2025}.\protect\footnote{KB excerpt:
\texttt{kb/excerpts/\glsterm{muon}{muon}-moonlight2025.md\#sec-1-challenges}.} The
mechanistic risk is that the polar-factor approximation in
\cref{eq:newton-schulz} breaks down when the spectrum of
$\mathbf{M}_t$ becomes too anisotropic at very large width, or when
distributed effects (gradient noise from large batches, \glsterm{expert-parallel}{EP} routing
variance) interact with orthogonalization in ways the small-scale
sweep does not exhibit. As of 2026-Q2 the result is single-source
on a 16B-MoE run; treat the $2\times$ claim as the strongest
\glsterm{llm}{LLM}-scale optimizer-efficiency evidence on record but not yet a
frontier-scale ground truth.
\end{contradiction}

\begin{contradiction}
\label{contra:training-loss-vs-downstream}
\glsterm{lion}{Lion} \citep[Chen et al.\ 2023]{muon-moonlight2025} and \glsterm{sophia}{Sophia}
\citep[Liu, Li et al.\ 2023]{muon-moonlight2025} both report
$\sim$1.5--2$\times$ training-loss speedups that have not held on
downstream evaluation in independent
\glsterm{llm}{LLM}-specific replications. \glsterm{per-shape-rms-scaling}{Muon}'s scaling-law claim
\eqref{eq:muon-scaling-fit} is also a training-loss claim ---
validation NLL on the pre-training distribution. The downstream
panel \citep[\S 3]{muon-moonlight2025}\protect\footnote{KB excerpt:
\texttt{kb/excerpts/\glsterm{muon}{muon}-moonlight2025.md\#sec-3-3-moonlight}.} of
the Moonlight tech report is supportive (Moonlight's downstream
metrics are reported competitive with similarly-sized \glsterm{precision-pinning}{DeepSeek-V3}
variants) but does not yet have an independent third-party
replication. The structural lesson stands: training-loss optimizer
gains must be re-verified on the downstream evaluation surface
before they translate into a deployment-relevant claim.
\end{contradiction}

Two adjacent open questions are folded into later sections.
Distributed \glsterm{per-shape-rms-scaling}{Muon}'s $32$k-GPU behavior and the $G = d \cdot t \cdot c
\cdot p \cdot e$ composition belong to \cref{sec:distributed-training}.
The \glsterm{per-shape-rms-scaling}{Muon}-update-RMS-vs-BF16-range coupling
\citep[\S 2.2]{muon-moonlight2025}, where vanilla
\cref{eq:muon-vanilla} without weight decay drives weight RMS past
BF16's high-precision range, is exactly the optimizer-times-precision
coupling \cref{sec:scaling-precision} and \cref{sec:mp-stability}
flag --- it reappears as the load-bearing reason
\cref{eq:muon-final}'s $0.2\sqrt{\max(A, B)}$ and weight-decay terms
are not optional. A third open question is the Schedule-Free
recipe's downstream effect on the WSD-cooldown-as-data-mixture-vehicle
role; the cooldown phase is where curated high-quality data is
loaded \eqref{eq:moonlight-schedule}, and removing the schedule
removes the natural data-mixture trigger.

\subsection{Forward link}
\label{sec:opt-forward}

This section fixed the optimizer (\glsterm{adamw}{AdamW} as the production default;
\glsterm{per-shape-rms-scaling}{Muon} as the only published 2024--2026 optimizer with a $\sim 2\times$
\glsterm{llm}{LLM}-scale compute-efficiency lift) and the schedule (WSD displacing
cosine, with Schedule-Free at the frontier). Two threads carry
forward. \Cref{sec:distributed-training} explores how each optimizer
composes with the full \glsterm{data-parallel}{DP} / \glsterm{tensor-parallel}{TP} / \glsterm{pipeline-parallel}{PP} / \glsterm{context-parallel}{CP} / \glsterm{expert-parallel}{EP} DeviceMesh ---
including the Distributed \glsterm{per-shape-rms-scaling}{Muon} \glsterm{data-parallel}{DP}-gather + Newton--Schulz partial
update of \cref{sec:opt-distributed-muon} as the latest add ---
because optimizer state and gradient communication are the two
bandwidth axes the parallelism plan must accommodate.
\Cref{sec:mixed-precision} explores how the \glsterm{adamw}{AdamW}
update-RMS calibration (\cref{tab:opt-shapes}, $0.2$--$0.4$) and
\glsterm{per-shape-rms-scaling}{Muon}'s matched $0.2\sqrt{\max(A, B)}$ scaling
(\cref{eq:muon-final}) interact with BF16's high-precision range,
why vanilla \cref{eq:muon-vanilla} requires the
\cref{eq:muon-wd}--\cref{eq:muon-final} sequence to stay numerically
stable, and how the format-zoo (FP32 master, BF16 working, \glsterm{fp8}{FP8}
GEMMs, \glsterm{nvfp4}{NVFP4} frontier) lands on the optimizer-precision boundary
this section opened. The next section's six-stage pre-training
scaffolding is, in effect, the engineering interface between this
section's optimizer and \cref{sec:scaling-laws}'s budget split.
